\section*{Chapitre \ref{ch:m2:settlement-risk} --- Risque de règlement}

\begin{solution}{exo:m2:settlement-risk:1}
Vendre des yens contre des dollars~: depuis l'heure limite des yens, à 22:00 la veille, jusqu'à la
confirmation du dollar à 21:00, soit 23 heures. Vendre des dollars contre des yens~: avec ces horaires
les yens arrivent avant que le paiement en dollars ne devienne irrévocable, si bien qu'il n'existe aucune fenêtre~;
en pratique, des heures limites d'instruction plus précoces peuvent en créer une.
\end{solution}

\begin{solution}{exo:m2:settlement-risk:2}
Brut~: la totalité du principal de la devise payée, puisque la jambe de la contrepartie
n'arrive jamais. Sous PvP~: seul le coût de remplacement de la transaction au nouveau cours, puisque
son propre paiement n'est pas effectué non plus.
\end{solution}

\begin{solution}{exo:m2:settlement-risk:3}
La première banque verse 100 millions d'euros nets, et les dollars sont compensés de la même manière~:
un seul paiement dans chaque devise au lieu de vingt.
\end{solution}

\begin{solution}{exo:m2:settlement-risk:4}
A~: la banque vend des euros contre 100 millions de dollars nets. B~: vend des yens contre 100 millions de dollars. C~:
vend des euros contre 100 millions de livres sterling, et des dollars contre 100 millions de livres sterling. D~:
vend la devise émergente pour 30 millions de dollars. 430 millions de dollars au total.
\end{solution}

\begin{solution}{exo:m2:settlement-risk:5}
Chaque jambe n'est réglée que lorsque les deux le sont, si bien que les membres doivent disposer de chaque devise qu'ils doivent
à CLS à temps, sans avoir d'abord reçu ce qui leur est dû dans l'autre~; il s'agit d'un
besoin de financement qui porte sur plusieurs devises et plusieurs horaires. Régler ensemble de nombreuses transactions permet
que les paiements entrants et sortants se compensent, de sorte que seuls les versements nets sont nécessaires.
\end{solution}

\begin{solution}{exo:m2:settlement-risk:6}
De nombreuses devises des marchés émergents ne sont pas réglées dans des systèmes PvP, leurs systèmes de paiement
peuvent ne pas se chevaucher avec les horaires du dollar, et les petites institutions n'ont souvent
pas d'accès direct ou abordable aux services PvP ou de \omterm{def:m2:settlement-risk:netting}{compensation des paiements}~: leurs transactions se règlent
en brut, par l'intermédiaire de comptes de correspondant.
\end{solution}

\begin{solution}{exo:m2:settlement-risk:7}
Brut~: 330, 930, 1\,330 et 730 millions de dollars à 05:00, 08:00, 14:00 et 17:00~;
après compensation~: 130, 330, 430 et 230 millions de dollars.
\end{solution}

\begin{solution}{exo:m2:settlement-risk:8}
La valeur de marché des forwards représente le coût de remplacement, soit l'exposition avant
le règlement. À la date de règlement, l'exposition est la totalité du principal de la
devise payée jusqu'à l'arrivée de l'autre, souvent des centaines de fois plus élevée, et elle
se reproduit à chaque date de valeur. Les expositions au règlement nécessitent leurs propres limites de règlement.
\end{solution}

\begin{solution}{pb:m2:settlement-risk:1}
\textbf{1.} 1.53~milliard de dollars d'achats~; un pic à 1.33~milliard de dollars.
\textbf{2.} De 13:00, moment où les paiements en dollars deviennent irrévocables, à 16:00, moment où
les euros et la livre sterling sont confirmés~: les ventes de dollars sont alors exposées, tandis que les
ventes antérieures d'euros et de yens attendent encore leurs dollars.
\textbf{3.} La vente de yens contre 250 millions de dollars à B~: 23 heures.
\textbf{4.} 330 millions de dollars~: la vente de yens par B et la vente de la devise émergente par D.
\textbf{5.} 700 millions de dollars~: les euros vendus contre des dollars (400 millions d'euros,
en attente des dollars) et les dollars vendus contre des euros (300 millions de dollars, en attente des euros).
\textbf{6.} Comme dans l'exercice 4~: 100, 100, 200 et 30 millions de dollars avec A, B, C et D.
\textbf{7.} 430 millions de dollars.
\textbf{8.} 100 millions de dollars.
\textbf{9.} Un accord de \omterm{def:m2:settlement-risk:netting}{compensation des paiements} opposable en cas d'insolvabilité de la contrepartie, de sorte
qu'un liquidateur ne puisse pas réclamer les montants bruts dus à la banque défaillante tout en ne
payant pas ceux qu'elle doit.
\textbf{10.} 30 millions de dollars~: la vente nette de la devise émergente, parce que cette
devise n'est pas réglée par PvP.
\textbf{11.} 30 millions de dollars.
\textbf{12.} Seul le risque de remplacement, une fraction du principal, puisque leurs transactions
sont réglées par \omterm{def:m2:settlement-risk:pvp}{paiement contre paiement}.
\textbf{13.} Les frais et les coûts d'adhésion ou d'accès, la liquidité nécessaire pour payer dans chaque
devise dans les temps, et les opérations nécessaires pour respecter les échéances.
\textbf{14.} Par un dispositif PvP pour cette devise, s'il en existe un, en demandant à D de payer
en premier, par un collatéral ou une limite, ou en réglant par l'intermédiaire d'une banque qui propose la
devise avec une protection contre les pertes.
\textbf{15.} Parce que la contrepartie n'a pas accès à PvP, la devise n'est pas
prise en charge, ou le coût est jugé trop élevé pour la transaction.
\textbf{16.} Sous la forme d'une limite de crédit sur l'exposition maximale au règlement par date de valeur,
dimensionnée en fonction de la qualité de crédit de la contrepartie et des fonds propres de la banque, et vérifiée
avant l'acceptation des transactions.
\textbf{17.} Il a montré que les transactions de change portent un risque de principal, et non seulement un risque de prix, et
a conduit à des limites de règlement, à la \omterm{def:m2:settlement-risk:netting}{compensation des paiements} et, à terme, à PvP.
\textbf{18.} Environ 1.4~trillion de dollars par jour reposent encore sur la confiance de chaque
payeur dans l'arrivée de l'autre devise~: une exposition importante et récurrente, concentrée
sur certaines devises et certaines contreparties.
\textbf{19.} Résultat nommé~: \emph{la journée non réglée}~: le pic d'exposition du livre est
1.33~milliard de dollars réglés en brut, 430 millions de dollars compensés, et 30 millions de dollars avec PvP
pour toutes les devises sauf une.
\textbf{20.} C'est le risque de payer une devise sans recevoir l'autre, et
le \omterm{def:m2:settlement-risk:pvp}{paiement contre paiement} le supprime.
\end{solution}

\begin{solution}{iq:m2:settlement-risk:1}
Le risque qu'une jambe d'une transaction de change soit payée de manière irrévocable, dans son propre fuseau horaire,
tandis que l'autre, échue plus tard ailleurs, n'est pas reçue parce que la contrepartie
fait défaut, comme lors de la fermeture de Herstatt en 1974. CLS règle les deux jambes dans ses livres au
même instant, ou aucune des deux, de sorte qu'aucun principal n'est exposé.
\iqlookfor{le mécanisme lié aux fuseaux horaires et PvP.}
\end{solution}

\begin{solution}{iq:m2:settlement-risk:2}
Le \omterm{def:m2:settlement-risk:risk}{risque de règlement} est la perte du principal entier lorsqu'une jambe est payée et que
l'autre n'arrive jamais. Le risque de remplacement est la perte liée à la nécessité de refaire une
transaction non réglée à un cours moins favorable après la défaillance de la contrepartie~: une fraction du principal,
liée aux mouvements de prix.
\iqlookfor{principal contre mouvement de marché.}
\end{solution}

\begin{solution}{iq:m2:settlement-risk:3}
Pour chaque transaction échue, déterminer l'instant où le paiement de la devise vendue devient irrévocable
et l'instant où la réception de la devise achetée est confirmée~; l'exposition à la
contrepartie à l'instant $t$ est la somme des montants achetés dont la date $t$ tombe dans cette fenêtre,
après compensation lorsque des accords le prévoient et hors transactions réglées en PvP~; compléter par
les réceptions non confirmées et échouées~; la calculer par date de valeur et en temps réel à mesure
que les confirmations arrivent.
\iqlookfor{fenêtres issues des échéanciers de paiement, compensation et confirmations.}
\end{solution}

\begin{solution}{iq:m2:settlement-risk:4}
Le PvP ne couvre que certaines devises et certaines plages horaires~; certaines contreparties n'y ont pas accès
directement ou par l'intermédiaire d'un membre~; y adhérer ou l'utiliser entraîne des frais et mobilise de la
liquidité~; certaines transactions (le jour même, tardives, certains produits) ne s'insèrent pas dans son
calendrier~; et certaines banques jugent le risque acceptable dans les limites fixées.
\iqlookfor{accès, couverture, coût et calendrier.}
\end{solution}

\begin{solution}{iq:m2:settlement-risk:5}
Dans la mesure où les flux des deux parties se compensent~: pour les teneurs de marché qui échangent
dans les deux sens toute la journée, la réduction est importante~; pour les clients qui n'opèrent que dans un
sens, elle est faible. Il faut des transactions de même date de valeur et de mêmes devises, un accord de
compensation juridiquement exécutoire, et des opérations qui conviennent des montants nets à temps.
\iqlookfor{compensation des flux et base juridique.}
\end{solution}

\begin{solution}{iq:m2:settlement-risk:6}
Tenir à jour, par contrepartie et par date de valeur, le profil d'exposition projeté à partir des
transactions enregistrées, des échéanciers de paiement et de la compensation~; à chaque nouvelle transaction,
calculer le profil en en tenant compte et comparer son pic à la limite avant acceptation, dans le
circuit de pré-négociation~; mettre à jour à mesure que les confirmations arrivent~; faire remonter ou
refuser lorsqu'une limite serait dépassée~; tester sur des journées historiques et des scénarios de défaillance.
\iqlookfor{contrôle préalable du pic projeté, mis à jour par les confirmations.}
\end{solution}
